\section{Conclusion}
\label{sec:conclusion}

We have described a marketplace in which the public catalog contains only what a seller chooses to publish, a buyer's shipping address is sealed to one seller and never enters the catalog, payments are non-custodial, and escrow is a two-party key in which a moderator can decide a dispute without ever being able to take the funds. Each mechanism carries a label. Signed catalogs, sealed orders, and the per-order wallet are implemented and tested. The two-party escrow is implemented and tested on an open pull request, but an internal review found open issues that block its use with real money. Tor delivery, a mailbox key bound to the store record, replica and seed networks, and private discovery are specified or future work.

The limits are as important as the mechanisms. The seller and the carrier must learn the ship-to. The chain records every payment. A network observer who can watch both ends can correlate timing. A moderator in a dispute learns the order. What the design changes is who else learns these things, and that no party can hold the buyer's or seller's funds alone. Before mainnet, the work that matters most is to fix the open review findings, deliver orders over Tor to a key bound to the store record, replace fixture keys with keys from the wallet seed, and obtain an external review of the escrow protocol.
